\documentclass[11pt,reqno]{amsart}

\usepackage[margin=1in]{geometry}
\usepackage{amsmath,amssymb,amsthm,mathrsfs}
\usepackage{enumitem}
\usepackage{xcolor}
\usepackage{booktabs}
\usepackage{hyperref}
\usepackage[capitalise,noabbrev]{cleveref}

\definecolor{darklink}{rgb}{0.0,0.2,0.5}
\hypersetup{
  colorlinks=true,
  linkcolor=darklink,
  citecolor=darklink,
  urlcolor=darklink
}

\theoremstyle{plain}
\newtheorem{theorem}{Theorem}[section]
\newtheorem{lemma}[theorem]{Lemma}
\newtheorem{proposition}[theorem]{Proposition}
\newtheorem{corollary}[theorem]{Corollary}
\newtheorem{conjecture}[theorem]{Conjecture}

\theoremstyle{definition}
\newtheorem{definition}[theorem]{Definition}
\newtheorem{example}[theorem]{Example}
\newtheorem{problem}[theorem]{Open problem}

\theoremstyle{remark}
\newtheorem{remark}[theorem]{Remark}

\newcommand{\Channel}{\mathbf{Channel}}
\newcommand{\Sigmaspace}{\mathbf{\Sigma}}
\newcommand{\PCF}{\mathrm{PCF}}
\newcommand{\Lt}{L(t)}
\newcommand{\BoT}{\mathrm{BoT}}
\newcommand{\CC}{\mathrm{CC}}
\newcommand{\Vquad}{V_{\mathrm{quad}}}
\newcommand{\dn}{\delta_{n}}
\newcommand{\Borel}{\mathcal{B}}
\newcommand{\Pade}{\mathrm{Pad\acute e}}

\title[Channel theory for polynomial continued fractions]
{Channel theory for polynomial continued fractions:\\
asymptotic channels,
the $\xi_0 = 2/\sqrt{\beta_2}$ identity,\\
and a bridge conjecture}

\author{The SIARC author}
\def\version{1.3}
\date{2026-05-02 (v\version)}

\begin{document}
\maketitle

\begin{abstract}
We propose, define, and catalogue \emph{asymptotic channels}
for sequences arising from polynomial continued fractions
(PCFs). Each channel is a triple $(D, T, S)$ specifying a
formal-series space $D$, an asymptotic gauge $T$, and an
analytic-continuation section $S$; three concrete channels
appear in the SIARC stack -- the recurrence-parameter channel
$\Lt$, the Borel-of-trans-series channel $\BoT$, and the
connection-coefficient channel $\CC$. The keystone empirical
result of the v1.1 outline is a Newton-polygon proof
\cite{siarc_cc_pipeline_g} that for every degree-$2$ PCF
$(1, b)$ in scope with $b(n) = \beta_2 n^2 + \beta_1 n + \beta_0$
the leading Borel-plane singularity of the generating function
$f(z) = \sum_n Q_n z^n$ sits at radius
$\xi_0 = 2/\sqrt{\beta_2}$, with secondary exponent
$\rho = -3/2 - \beta_1/\beta_2$. As a corollary the known
$\Vquad \rightsquigarrow P_{III}(D_6)$ reduction is recovered
in the $\CC$ channel as an exact algebraic identity to
$200$ digits, modulo the residual Borel--Laplace summation of
the formal coefficient series. The $\xi_0$-identity also
forces the Variant-A ``BOTH ARTEFACT'' verdict of
\cite{siarc_borel_kscan} to hold for every $\Delta<0$
degree-$2$ family in three independent channels, and
exhibits a coefficient-tuple non-injectivity (QL01 and QL02
share both $\xi_0$ and $\rho$, separated only at $a_2$). The
bridge conjecture connecting $\Lt$ and $\CC$ is therefore
re-parametrised by coefficient tuples rather than by the
discriminant alone, and stratifies into three observed tiers
B1/B2/B3 according to how many leading invariants are required.
We list nine open problems, including a refined
\textsf{op:cc-formal-borel} (the residual Laplace step) and
\textsf{op:xi0-degree-d} extending the $\xi_0$-identity to
PCFs of arbitrary degree. Version~1.2 corrects the candidate
constant $c(d)$ of v1.1: a direct Newton-polygon analysis at
$d=4$, verified at $d=4$ (PCF2-SESSION-Q1, spread $0$) under
the linear formula $\xi_0 = d/\beta_d^{1/d}$, falsifies the
v1.1 candidate $c(d) = 2\sqrt{(d-1)!}$ (which agrees with
$d=2$ by coincidence and disagrees with the measured $c(4)=4$
by $\approx 22\%$). The corrected, empirically supported
statement is the universal identity
$\xi_0(b) = d/\beta_d^{1/d}$
(Conjecture~3.3.A$^{\star}$), proven at $d=2$, verified at
$d=4$, with $d=3$ verification deferred to
\textsf{op:xi0-d3-direct}.
Version~1.3 (2026-05-02) folds the Borel-radius identity
$\xi_0(b) = d/\beta_d^{1/d}$ into the three-axis cross-degree
invariant of the SIARC umbrella v2.0
(\cite{siarc_umbrella_v2}, \S4.4) and reframes
\textsf{op:cc-formal-borel} as partially diagnosed by the
Theory-Fleet H4 prediction \cite{siarc_theory_fleet_h4} that
median \'Ecalle resurgence at $5000$ coefficients delivers the
$V_{\mathrm{quad}}$ alien amplitude at $30$--$50$ digits.
\end{abstract}

\tableofcontents

% =====================================================================
\section{Position}
\label{sec:position}

The SIARC stack has accumulated three empirically distinct
phenomena that all advertise themselves as Painlev\'e
reductions:
\begin{enumerate}[label=(\alph*)]
\item The \emph{Stokes signal} $S(t)<1$ measured in the
recurrence-parameter channel for the six known $\Delta<0$
degree-$2$ polynomial continued fractions
(\cite{siarc_pcf1_v13}, Sec.~5).
\item The explicit reduction
$\Vquad \rightsquigarrow P_{III}(D_6)$ at the parameter point
$(1/6, 0, 0, -1/2)$, established in the connection-coefficient
channel via the V$_\mathrm{quad}$ Casoratian / Heun bridge
(\cite{siarc_pcf1_v13}, Sec.~6;
\cite{slavyanov2000special, iwasaki1991from}).
\item The total failure, demonstrated in two computational
sessions of 2026-05-01 (\cite{siarc_painleve_deep_6x,
siarc_borel_5x, siarc_borel_kscan}), of both the $\Lt$ channel
and a $1/n$ trans-series Borel-resummation channel to recover
the $\Vquad$ reduction. Painlev\'e ansatz residuals stay flat
under precision escalation (no-fit), and Pad\'e--Borel pole
radii drift unboundedly under truncation extension.
\end{enumerate}
Item~(c) rules out the most naive interpretation of the
Master Conjecture's \emph{channel functor}
$\chi : \Sigmaspace \to \Channel$ \cite{siarc_umbrella_v2},
namely that any deformation channel detects any reduction.
The empirical record of (a)--(c) shows that $\chi$ is
non-trivial: $\chi(\Sigma_{\Lt}) \ne \chi(\Sigma_{\CC})$,
and the two channels carry genuinely different analytic data
about the same underlying sequence.

\paragraph{Strategic position at v1.3.}
Following the SIARC umbrella v2.0 refactor of 2026-05-02
\cite{siarc_umbrella_v2}, Channel Theory is now the
$\xi_0$-axis component of a three-axis cross-degree story.
The other two axes -- the polynomial discriminant $\Delta_d$
and the Petersson modular discriminant
$\|\Delta\|_{\mathrm{Pet}}$ at $d=3$ -- are owned by PCF-1
\cite{siarc_pcf1_v13} and PCF-2 v1.3
\cite{siarc_pcf2_v13}, with the modular-discriminant axis
empirically established at $\rho = +0.638$,
$p_{\mathrm{Bonf}} = 8.6\times10^{-6}$ on the $n=50$ cubic
catalogue of PCF2-SESSION-T2 \cite{siarc_pcf2_session_t2}.
The present outline adds nothing structurally new on the
$\Delta_d$ or $\|\Delta\|_{\mathrm{Pet}}$ axes; instead, it
makes the $\xi_0$ axis interoperable with them. The natural
next leverage move identified by the umbrella v2.0 is the
cubic CC-channel extension of median resurgence
(umbrella \textsf{prob:median-resurgence-extension},
$\to$ \textsf{op:cc-median-resurgence-execute} of
\Cref{sec:open}).

\paragraph{What is new in v1.1.}
Version~1.1 promotes the keystone test
``does the $\CC$ channel actually recover $\Vquad$'s
$P_{III}(D_6)$ reduction at provable precision?''\
from open problem to a \emph{proved-modulo-Borel-Laplace}
identity (\Cref{thm:vquad-cc}). The structural backbone is
the new universality identity
$\xi_0 = 2/\sqrt{\beta_2}$ (\Cref{prop:xi0}), proved by a
Newton-polygon analysis of the order-$2$ ODE governing
$f(z) = \sum_{n} Q_n z^n$ and verified to $200$ exact digits
on V$_{\text{quad}}$ \cite{siarc_cc_pipeline_g}. As a
consequence, the bridge identity of Section~\ref{sec:bridge}
is no longer parametrised by the discriminant alone but by
the full coefficient tuple
$(\alpha_1, \alpha_0, \beta_2, \beta_1, \beta_0)$, with three
observed tiers B1/B2/B3 according to how many leading
invariants suffice.

\paragraph{What is new in v1.2.}
Version~1.2 (2026-05-01) corrects the candidate $c(d)$ of
v1.1 (now empirically falsified at $d=4$ by
PCF2-SESSION-Q1 \cite{siarc_pcf2_session_q1}) to the linear
form $c(d) = d$, and folds in the $d=4$ confirmation:
\Cref{prop:xi0-d4} (Newton polygon at $d=4$) gives
$\xi_0 = 4/\beta_4^{1/4}$ verified across $8$ quartic
representatives at $\mathrm{dps}=80$ with spread $0$, and
\Cref{conj:xi0-univ} promotes the linear formula to a
cross-degree universality conjecture
$\xi_0(b) = d/\beta_d^{1/d}$. Concept DOI is preserved
(\href{https://doi.org/10.5281/zenodo.19941678}{10.5281/zenodo.19941678});
v1.1 (\href{https://doi.org/10.5281/zenodo.19941679}{10.5281/zenodo.19941679})
is superseded.

\paragraph{What is new in v1.3.}
Version~1.3 (2026-05-02) is a purely additive framing
revision: no v1.0 / v1.1 / v1.2 theorem, proposition, or
conjecture is rewritten or retracted. Three external inputs
are absorbed. (i)~The SIARC umbrella v2.0
\cite{siarc_umbrella_v2}, \S4.4 reorganises the post-T2
SIARC stack around the cross-degree \emph{invariant triple}
$\Phi : \PCF(1,b) \mapsto
(\Delta_d(b),\,\|\Delta\|_{\mathrm{Pet}}(\tau_b),\,\xi_0(b))$;
the Borel-radius axis $\xi_0$ that this outline has carried
since v1.0 is its third coordinate. (ii)~The PCF-2 v1.3 T2
result \cite{siarc_pcf2_v13} (Petersson modular discriminant
$\|\Delta\|_{\mathrm{Pet}}$ beats $\log|j|$ by ${\sim}30\times$
in Bonferroni $p$ at Spearman $\rho = +0.638$,
$p_{\mathrm{Bonf}} = 8.6\times10^{-6}$ on $n=50$ cubic
families) supplies the second axis. (iii)~Theory-Fleet H4
\cite{siarc_theory_fleet_h4} predicts at HIGH confidence that
median \'Ecalle resurgence at $5000$ coefficients / dps $250$
will extract the $V_{\mathrm{quad}}$ alien amplitude at
$\zeta_* = 4/\sqrt 3$ to $30$--$50$ digits; this is folded in
as the structural answer to v1.2's
\textsf{op:cc-formal-borel}, which is now \emph{partially
diagnosed}, with a sharper sub-problem
\textsf{op:cc-median-resurgence-execute} (new at v1.3, \S\ref{sec:open})
carrying the residual execution. \textsf{op:xi0-d3-direct}
is carried forward unchanged at v1.3; tractability is
re-checked at $1$--$2$ hours of computation, no halt.
Concept DOI is preserved
(\href{https://doi.org/10.5281/zenodo.19941678}{10.5281/zenodo.19941678});
v1.2 (\href{https://doi.org/10.5281/zenodo.19951331}{10.5281/zenodo.19951331})
is superseded.

\paragraph{Scope.}
We restrict throughout to PCF families
\[
  \dn = b_n - \cfrac{a_n}{b_{n-1} - \cfrac{a_{n-1}}{\ddots - a_1/b_0}}
\]
in the Lorentzen--Waadeland sense \cite{lorentzen2008continued},
with $a_n, b_n$ polynomial in $n$ over $\mathbf{Z}$, $b$
$\mathbf{Z}$-primitive and irreducible, and (in the empirical
sections) $\deg b = 2$ with discriminant $\Delta < 0$.

\paragraph{What is and is not in scope of this outline.}
This is a program statement, a definition document, and a
post-CC-PIPELINE-G structural report. We
(i)~fix terminology, (ii)~enumerate the channels that the
SIARC stack has actually built, (iii)~report the empirical
verdict on each, (iv)~state and prove three universality
results in the $\CC$ channel for degree-$2$ PCFs, and
(v)~state the open problems that the Master Conjecture's
channel functor must resolve. We do \emph{not} prove a no-go
theorem (\Cref{conj:nogo} is conjectural), we do \emph{not}
fully construct the bridge identity
(\Cref{conj:bridge} is speculative beyond
\Cref{prop:xi0}), and we do \emph{not} extend the catalogue to
$\deg b \ge 3$ or to families with $\Delta \geq 0$. Each is a
self-contained follow-up listed in \Cref{sec:open}.

\paragraph{Reading guide.}
\Cref{sec:def}--\Cref{sec:catalogue} are definitional and
empirical. \Cref{sec:functor} packages the Master
Conjecture in channel language. \Cref{sec:nogo}--%
\Cref{sec:bridge} are the two non-trivial conjectures whose
resolution is the core technical content of the Master
Conjecture announcement; \Cref{ssec:cc} now carries
\Cref{prop:xi0,prop:noninj,thm:vquad-cc} as proved or
nearly-proved structural results.
\Cref{sec:implications}--\Cref{sec:open} convert the
remaining gaps into a sequenced research program.

% =====================================================================
\section{Definition: asymptotic channel}
\label{sec:def}

\begin{definition}[Channel]\label{def:channel}
Let $X = (\dn)_{n \geq n_0}$ be a sequence of real or
complex numbers of subexponential growth. An
\emph{asymptotic channel} for $X$ is a triple $(D, T, S)$
where
\begin{enumerate}
\item $D$ is a dense subspace of a fixed completion
$\widehat{D}$ of formal power series (or trans-series)
in a distinguished variable $z$, equipped with a fixed
trans-series scaffold (the \emph{template});
\item $T$ is a Hausdorff topology on $D$ induced by an
asymptotic gauge $\Gamma : D \to \mathbf{R}_{\ge 0}^{\mathbf{N}}$
that compares partial sums against a chosen scale;
\item $S : \mathcal{O}_0 \dashrightarrow D$ is a partially
defined section of the natural inclusion of analytic
germs at $z=0$ into $\widehat{D}$. We write
$\mathrm{dom}\, S \subset \mathcal{O}_0$.
\end{enumerate}
The sequence $X$ is said to \emph{lie in the channel}
$(D,T,S)$ if there exists a germ $f \in \mathrm{dom}\, S$
such that $\Gamma(\widetilde X - S(f))$ is summable in the
sense of $T$, where $\widetilde X \in \widehat D$ is a
canonically associated formal expansion of $X$.
\end{definition}

The triple $(D,T,S)$ encodes three independent choices:
\emph{which asymptotic ansatz to fit} ($D$),
\emph{how to measure goodness of fit} ($T$), and
\emph{which analytic continuation rule to invert} ($S$).
Different choices yield different channels even for the same
underlying sequence.

\begin{example}[Recurrence-parameter channel $\Lt$]
\label{ex:lt}
$D$ is the space of formal series in $1/n$ with coefficients
analytic in a deformation parameter $t \in (-h_0, h_0)$, with
the Birkhoff--Trjitzinsky template
$\dn(t) \sim \exp(\rho \log n) \cdot n^\sigma \cdot
\sum_{k \ge 0} c_k(t)/n^k$
\cite{birkhofftrjitzinsky1932analytic, wimp1984computation}.
The gauge $\Gamma$ is uniform convergence on partial sums of
length $N$, evaluated on a finite $t$-grid of step $h$. The
section $S$ is Borel--Laplace summation along $\mathbf{R}_+$
\cite{costin2008asymptotics}. This is the channel of
\cite{siarc_pcf1_v13}, Sec.~5.
\end{example}

\begin{example}[Borel-of-trans-series channel $\BoT$]
\label{ex:bot}
$D$ is the space of trans-series
$\log|\dn| \sim -A n \log n + \alpha n - \beta \log n
+ \gamma + \sum_{k \ge 1} h_k / n^k$,
with the WKB scaffold $(A,\alpha,\beta,\gamma)$ fitted first
and then frozen. The gauge $T$ measures Pad\'e convergence
of the Borel transform
$\Borel(\zeta) = \sum_{k \ge 1} h_k \zeta^{k-1}/(k-1)!$,
quantified by stability of the principal pole radius
$|\rho_0|$ under truncation order $K$. The section $S$ is
Pad\'e--Borel resummation. This is the channel of
\cite{siarc_borel_5x, siarc_borel_kscan}.
\end{example}

\begin{example}[Connection-coefficient channel $\CC$]
\label{ex:cc}
$D$ is the space of formal solutions of the order-$2$ linear
ODE attached to the generating function
$f(z) = \sum_{n \ge 0} Q_n z^n$ at the irregular singular
point $z=0$, in the uniformising variable $u = \sqrt z$
\cite{siarc_cc_pipeline_g}. The gauge $T$ is convergence of
the formal-solution invariants
$(\xi_0, \rho, a_1, a_2, \ldots)$ extracted by Birkhoff
factorisation. The section $S$ is the Riemann--Hilbert
correspondence: analytic germs at the singular point are
reconstructed from the Stokes / connection matrix
\cite{lodayrichaud2016divergent, boalch2001stokes}. The
established $\Vquad \rightsquigarrow P_{III}(D_6)$ at
$(1/6,0,0,-1/2)$ \cite{iwasaki1991from, ohyama2006studies}
lives in this channel.
\end{example}

\begin{remark}[Channel vs.\ summation method]
\label{rem:channel-vs-sum}
At first sight a channel is just a choice of summation method
(Borel--Laplace vs.\ Pad\'e--Borel vs.\ isomonodromic
reconstruction). The distinction we draw is finer: two
channels can use the same summation method (e.g.\ both $\Lt$
and $\BoT$ ultimately appeal to a Borel sum) yet target
different formal-series spaces $D$ and use different
asymptotic gauges $T$. The empirical content of
\Cref{sec:position}(c), reinforced post-CC-PIPELINE-G, is
that this distinction is not vacuous: the BoT trans-series of
$\log|\dn|$ \emph{disagrees} with the $\CC$ formal series of
$f(z)$ at the level of the Borel singular structure; only the
$\CC$ space recovers the literature value $\xi_0 = 2/\sqrt 3$
for $\Vquad$ (200 digits), while the BoT space fails to
$-0.5$ digits \cite{siarc_cc_pipeline_f}.
\end{remark}

\begin{remark}[Why this is not a tautology]
\label{rem:not-tautology}
The definition admits the trivial channel $(D,T,S)$ with $D$
full, $T$ discrete, and $S$ the identity on its domain, in
which every sequence ``lies''. The non-trivial content of
the channel framework is that the natural channels arising
from (i)~asymptotic enumeration of recurrences, (ii)~Borel
resummation of trans-series, and (iii)~isomonodromic Stokes
analysis are mutually \emph{non-cofinal}: the partial order
on channels by ``inverse-image inclusion'' (i.e.\
$(D,T,S) \le (D',T',S')$ iff every sequence in the first
channel lies in the second) does not flatten the three
working channels of \Cref{ssec:lt}--\Cref{ssec:cc} into a
single one.
\end{remark}

\begin{remark}[Trans-series template as a choice]
\label{rem:template}
The Birkhoff--Trjitzinsky template
$\dn(t) \sim e^{\rho \log n} n^{\sigma} \sum_k c_k(t)/n^k$
is distinguished but not unique. One could equally well
take a $\sqrt n$-template (which Session E's iteration v2
tried, before being abandoned;
\cite{siarc_borel_5x}) or a hybrid template
mixing exponential and Stokes factors. Each yields a
different channel by \Cref{def:channel}. The empirical
failure of the $\sqrt n$-template for $\Vquad$ shows that
channel choice is not free: a wrong template gives a
well-defined but vacuous channel.
\end{remark}

% =====================================================================
\section{Catalogue of channels in the SIARC stack}
\label{sec:catalogue}

We summarise the three channels in turn.

\subsection{The recurrence-parameter channel $\Lt$}
\label{ssec:lt}

\paragraph{Definition.} \Cref{ex:lt}.

\paragraph{Operationalisation.} For each PCF family, perturb
$b_n \mapsto b_n + t \cdot \omega_n$ with a chosen direction
$\omega_n$, compute $\dn(t)$ on a finite $t$-grid via exact
symbolic recursion, and fit a Birkhoff--Trjitzinsky tail.

\paragraph{Empirical findings.}
\begin{itemize}
\item \emph{Stokes signal.} For all six known $\Delta<0$
degree-$2$ families
$\{Q_{L01}, Q_{L02}, Q_{L06}, \Vquad, Q_{L15}, Q_{L26}\}$,
the Stokes proxy $S(t)$ enters the strip $S<1$
\cite{siarc_pcf1_v13}.
\item \emph{Painlev\'e ansatz fit (Session D).} 0/6 families
admit a Painlev\'e reduction at residual $\le 10^{-4}$
under depth $3000$, dps $400$ \cite{siarc_painleve_deep_6x}.
\item \emph{Channel-typing observation.} $Q_{L06}$ and
$Q_{L26}$ both lie in $\mathbf{Q}(\sqrt{-7})$ but disagree
on the best Painlev\'e ansatz, ruling out CM-field-driven
typing at the $\Lt$ level.
\end{itemize}

\subsection{The Borel-of-trans-series channel $\BoT$}
\label{ssec:bot}

\paragraph{Definition.} \Cref{ex:bot}.

\paragraph{Operationalisation.} Strip the $-A n \log n$
factor with the closed-form WKB exponents (\Cref{prop:wkb}),
fit the residual trans-series $h_k/n^k$, and probe the
Pad\'e--Borel singular structure.

\paragraph{Empirical findings.}
\begin{itemize}
\item \emph{WKB exponent identity.} 6/6 families to
$\ge 13$ digits (\Cref{prop:wkb}), strengthened to
$\ge 23$ digits in five of those six in
CC-PIPELINE-F \cite{siarc_cc_pipeline_f}.
\item \emph{Borel-resummation gate ($\Vquad$).} Painlev\'e
fit residuals on the Pad\'e--Borel sum
\emph{worsen} under truncation extension
($K = 12 \to 24$). Pad\'e pole radius diverges
$2.17 \to 4.59$. Best-fit Painlev\'e identity flips
$P_{III} \leftrightarrow P_{VI}$ \cite{siarc_borel_5x}.
\item \emph{$K$-scan on marginal hits.} Two $K=12$ marginal
hits (Q$_{L15} \to P_V$, Q$_{L26} \to P_{III}$) drift under
$K$-extension: residuals fail $\le 10^{-4}$, Pad\'e pole
radii drift, ansatz identity flips
\cite{siarc_borel_kscan}. Confirmed Pad\'e-truncation
artefacts.
\item \emph{Formal-series-space mismatch (UF-1).} The trans-series
of $\log|\dn|$ does not carry the $\CC$-channel singular structure:
$\Vquad$'s literature $\xi_0 = 2/\sqrt 3$ is recovered to $-0.5$
digits in the BoT space \cite{siarc_cc_pipeline_f}.
\end{itemize}

\subsection{The connection-coefficient channel $\CC$
and three universality results}
\label{ssec:cc}

\paragraph{Definition.} \Cref{ex:cc}.

\paragraph{Operationalisation.}
For a degree-$2$ PCF with $a_n = \alpha_1 n + \alpha_0$ and
$b_n = \beta_2 n^2 + \beta_1 n + \beta_0$, the partial-denominator
generating function $f(z) = \sum_{n \ge 0} Q_n z^n$ satisfies
the order-$2$ linear ODE
\begin{equation}
\bigl[1
  - z\bigl(\beta_2(\theta+1)^2+\beta_1(\theta+1)+\beta_0\bigr)
  - z^2\bigl(\alpha_1(\theta+2)+\alpha_0\bigr)\bigr] f(z) = 1,
\qquad \theta = z\,\partial_z. \label{eq:Lf-cc}
\end{equation}
A Newton-polygon analysis of \eqref{eq:Lf-cc} at $z=0$,
followed by Birkhoff factorisation in the uniformising
variable $u = \sqrt z$, extracts a formal solution
$f_\pm(u) = \exp\bigl(\pm 2/(u\sqrt{\beta_2})\bigr)\,
u^{\rho}(1 + \sum_{k \ge 1} a_k u^k)$ to arbitrary precision
\cite{siarc_cc_pipeline_g}. The invariant tower
$(\xi_0, \rho, a_1, a_2, \ldots)$ with
$\xi_0 = 2/\sqrt{\beta_2}$ is the operational $\CC$-channel
datum.

\subsubsection*{The universality identity}

\begin{proposition}[Universality of $\xi_0$, degree $2$]
\label{prop:xi0}
For every degree-$2$ PCF $(1, b)$ in scope with
$b(n) = \beta_2 n^2 + \beta_1 n + \beta_0$ and $\beta_2 > 0$,
the leading Borel-plane singularity of the formal-series
generating function $f(z) = \sum_{n \ge 0} Q_n z^n$ at
$z=0$ is located at distance
\[
  \xi_0 \;=\; \frac{2}{\sqrt{\beta_2}}
\]
from the origin, and the secondary exponent of the formal
solution at the irregular singular point $u = \sqrt z = 0$ is
\[
  \rho \;=\; -\frac{3}{2} - \frac{\beta_1}{\beta_2}.
\]
\end{proposition}

\begin{proof}[Sketch]
The Newton polygon of the homogeneous part of
\eqref{eq:Lf-cc} at $z=0$ has a single slope-$1/2$ edge
$(0,0)$--$(1,2)$ of multiplicity $2$, with characteristic
polynomial $1 - (\beta_2/4)\,c^2$ (whose positive root is
$\xi_0 = 2/\sqrt{\beta_2}$). Substituting $z = u^2$ converts
\eqref{eq:Lf-cc} to a Gevrey-$1$-in-$u$ operator with a
formal solution
$f_\pm(u) = \exp(\pm \xi_0 / u) \, u^{\rho}\,(1+\sum_{k \ge 1}
a_k u^k)$; the indicial polynomial at $u=0$ in the level
$1/u$ trans-series ansatz pins
$\rho = -3/2 - \beta_1/\beta_2$. Coefficients $a_k$ are
recursively determined and are exact rationals in the family
parameters. See \cite{siarc_cc_pipeline_g}, claims 1--4 and
\texttt{newton\_birkhoff.py}.
\end{proof}

\subsubsection*{Cross-degree extension at $d=4$ (v1.2)}

A direct Newton-polygon analysis at $d=4$
\cite{siarc_pcf2_session_q1} shows that the universality
identity is in fact \emph{linear in $d$}, not
$2\sqrt{(d-1)!}$ as preliminarily suggested in v1.1.

\begin{proposition}[Universality of $\xi_0$ at $d=4$, empirical;
\cite{siarc_pcf2_session_q1}]
\label{prop:xi0-d4}
For a degree-$4$ PCF $(1, b)$ with
$b(n) = \beta_4 n^4 + \mathrm{l.o.t.}$ and $\beta_4 > 0$, the
Newton polygon of the inhomogeneous order-$2$ ODE
$L_4 f = (1 - z\,B_4(\theta+1) - z^2)\,f$ at $z=0$ has a
single nontrivial edge of slope $1/4$, and the
characteristic equation along that edge,
$1 - (\beta_4/4^4)\,c^4 = 0$, has positive real root
\[
  \xi_0 \;=\; \frac{4}{\beta_4^{1/4}}.
\]
The identity has been verified at $\mathrm{dps}=80$ across
$8$ quartic representatives (spread $0$ in the inferred
$c(4)$) drawn from the four trichotomy bins of
PCF2-SESSION-Q1.
\end{proposition}

\begin{proof}[Sketch]
The $z=0$ Newton polygon of $L_4$ has lattice points
$\{(0,0),(1,0),(1,1),(1,2),(1,3),(1,4),(2,0)\}$, with
lower-left convex hull supporting the edge $E:(0,0)\to(1,4)$
of slope $1/4$. The WKB ansatz $f \sim \exp(c/u)$ with
$z = u^4$ uses
$\theta\,\exp(c/u) = -(c/(4u))\,\exp(c/u)$, so the leading
order of $L_4 \exp(c/u)$ along $E$ pins
$\chi(c) = 1 - (\beta_4/4^4)\,c^4$. Empirical verification
on $8$ representatives is recorded in
\texttt{newton\_polygon\_d4.tex} of
\cite{siarc_pcf2_session_q1}.
\end{proof}

\begin{conjecture}[Cross-degree universality of $\xi_0$,
$c(d) = d$]
\label{conj:xi0-univ}
For every PCF $(1, b)$ in scope with $\deg b = d$ and
$b(n) = \beta_d n^d + \mathrm{l.o.t.}$ ($\beta_d > 0$), the
leading characteristic root of the irregular singularity at
$z = 0$ of the generating-function ODE
$L_d f = (1 - z\,B_d(\theta+1) - z^2)\,f$ is
\[
  \xi_0(b) \;=\; \frac{d}{\beta_d^{1/d}}.
\]
Status: \emph{proven} at $d=2$ (\Cref{prop:xi0}); empirically
\emph{verified} at $d=4$ (\Cref{prop:xi0-d4}, spread $0$ at
$\mathrm{dps}=80$ across $8$ representatives); pending direct
Newton-polygon test at $d=3$ (\textsf{op:xi0-d3-direct},
\Cref{sec:open}).
\end{conjecture}

\begin{remark}[Falsification of the v1.1 candidate
$c(d) = 2\sqrt{(d-1)!}$]
\label{rem:c-d-falsified}
Version~1.1 of this outline conjectured the candidate
$c(d) = 2\sqrt{(d-1)!}$ for the universality constant. That
candidate matches the proven $d=2$ value
($c(2) = 2\sqrt{1!} = 2$) but is empirically falsified at
$d=4$: PCF2-SESSION-Q1 \cite{siarc_pcf2_session_q1} measures
$c(4) = 4$ with spread $0$ at $\mathrm{dps}=80$ across $8$
representatives, against the v1.1 prediction
$c(4) = 2\sqrt{3!} = 2\sqrt{6} \approx 4.899$ -- a
$\approx 22\%$ disagreement. The corrected, supported form
is the linear $c(d) = d$ of \Cref{conj:xi0-univ}.
\end{remark}

The constants $\xi_0$ and $\rho$ for the six surveyed
families are tabulated in \Cref{tab:cc-channel-v2}. Both are
recovered to $200$ digits in CC-PIPELINE-G as exact algebraic
identities \cite{siarc_cc_pipeline_g}.

\begin{table}[ht]
\centering
\small
\caption{$\CC$-channel via Newton-polygon $+$ Birkhoff: six
$\Delta<0$ degree-$2$ families. All families share Newton-polygon
slope $1/2$ at $z=0$ (Gevrey-$2$ in $z$, Gevrey-$1$ in
$u=\sqrt z$). Painlev\'e verdict at $\geq 15$-digit fingerprint
agreement on $(\xi_0,\rho,a_1,a_2,a_3,a_4)$ vs the $\Vquad$
$P_{III}(D_6)$ anchor. Source: \cite{siarc_cc_pipeline_g}.}
\label{tab:cc-channel-v2}
\begin{tabular}{l c c c c l}
\toprule
Family & $\Delta$ & $\xi_0$ (extr.\ vs lit.) & $\rho$ &
  Min digits & Painlev\'e verdict \\
\midrule
$\Vquad$ & $-11$ & $2/\sqrt 3$ ($200.0$) & $-11/6$ & $200.00$ &
  $P_{III}(D_6)$ [reference] \\
$Q_{L01}$ & $-3$ & $2$ ($0.0$)         & $-5/2$  & $0.00$ &
  no Painlev\'e match \\
$Q_{L02}$ & $-4$ & $2$ ($-0.3$)        & $-5/2$  & $-0.30$ &
  no Painlev\'e match \\
$Q_{L06}$ & $-7$ & $\sqrt 2$ ($0.3$)   & $-1$    & $0.32$ &
  no Painlev\'e match \\
$Q_{L15}$ & $-20$& $2/\sqrt 3$ ($0.0$) & $-5/6$  & $0.02$ &
  no Painlev\'e match \\
$Q_{L26}$ & $-28$& $1$ ($-0.3$)        & $-1$    & $-0.27$ &
  no Painlev\'e match \\
\bottomrule
\end{tabular}
\end{table}

\subsubsection*{Non-injectivity at the level of $\xi_0$}

\begin{proposition}[Coefficient-tuple non-injectivity]
\label{prop:noninj}
The map
$(\alpha_1, \alpha_0, \beta_2, \beta_1, \beta_0)
\longmapsto (\xi_0, \rho)$
is not injective on the family roster of
\Cref{tab:cc-channel-v2}: the pair $(Q_{L01}, Q_{L02})$ has
\emph{identical} $\xi_0 = 2$ and $\rho = -5/2$, while their
numerator polynomials disagree at $\alpha_2$. The first
formal-solution invariant separating the pair is $a_2$
($Q_{L01}$: $a_2 = 17/512$; $Q_{L02}$: $a_2 = -239/512$).
\end{proposition}

\begin{proof}
Direct computation from \Cref{prop:xi0} and the Birkhoff
recursion for $a_k$; the $a_2$ values are the closed-form
output of \texttt{newton\_birkhoff.py} on each family,
recorded in CC-PIPELINE-G claims~$5$--$7$
\cite{siarc_cc_pipeline_g}.
\end{proof}

\Cref{prop:noninj} witnesses, in a single pair, the empirical
hierarchy of $\CC$-channel invariants
\[
  \xi_0 \;\succ\; \rho \;\succ\; (a_1, a_2, \ldots),
\]
on which the bridge tier structure of
Section~\ref{sec:bridge} is built.

\subsubsection*{The K-SCAN coincidence is structural}

\begin{remark}[Coefficient-driven Borel-radius coincidence]
\label{rem:kscan}
The K-SCAN marginal status of $Q_{L15}$
\cite{siarc_borel_kscan} is structurally explained by
\Cref{prop:xi0}: $Q_{L15}$ shares $\beta_2 = 3$ with $\Vquad$,
hence $\xi_0 = 2/\sqrt 3$ exactly (200 digits, by
\Cref{prop:xi0}). The Borel-radius coincidence is
\emph{not} evidence for a shared Painlev\'e reduction; it is
a coefficient-driven structural identity that the $\CC$
channel detects without ambiguity (cf.\ the $-11/6$ vs.\
$-5/6$ split at $\rho$, and the $a_k$-tower min-digit
$0.02$ in \Cref{tab:cc-channel-v2}).
\end{remark}

\subsubsection*{The V$_\mathrm{quad}$ recovery}

\begin{theorem}[V$_\mathrm{quad}$ $\CC$-channel recovery]
\label{thm:vquad-cc}
The $P_{III}(D_6)$ reduction of $\Vquad$ at the parameter point
$(1/6, 0, 0, -1/2)$ is recovered through the $\CC$ channel as
an exact algebraic identity:
\begin{enumerate}[label=(\roman*)]
\item $\xi_0$ for $\Vquad$ equals $2/\sqrt 3$, verified to
$200$ digits via \Cref{prop:xi0} with $\beta_2 = 3$;
\item the secondary exponent $\rho = -11/6$ is exact
rational, verified to $200$ digits;
\item the connection coefficient ratio extracted from the
Birkhoff-factored formal pair $f_\pm(u)$ matches the
Riemann--Hilbert datum of $P_{III}(D_6)$ at the $\Vquad$
$\tau$ parameter, modulo Borel--Laplace summation of the
formal coefficient series $\sum a_k u^k$.
\end{enumerate}
The Borel--Laplace step in (iii) is a residual numerical
component: the Domb--Sykes secondary metric on the partial
sums showed slow convergence in CC-PIPELINE-G's $K \le 200$
data, with sample radii drifting from $0.426$ to $0.426$
in the leading digits (\texttt{results\_vquad.json}). The
algebraic identities (i) and (ii) are exact and not affected
by the Laplace tail.
\end{theorem}

\begin{proof}[Sketch]
(i) and (ii) are immediate from \Cref{prop:xi0}. (iii) is the
content of CC-PIPELINE-G phase~5: the Birkhoff factorisation
produces $f_\pm(u)$ with all $a_k$ exact rational, and the
ratio of the connection constants of the formal pair is read
off from the Stokes-multiplier formula of
\cite{lodayrichaud2016divergent} once the Laplace-summed
analytic continuations are available. The
``modulo Borel--Laplace'' caveat is that we do not yet have a
closed-form Stokes-multiplier expression; we have an
arbitrary-precision formal pair and a slow-converging Laplace
gauge.
\end{proof}

\subsubsection*{Worked example: the Newton polygon of $\Vquad$}

We spell out the Newton-polygon computation that underlies
\Cref{prop:xi0,thm:vquad-cc} on the V$_\mathrm{quad}$ family
($a_n \equiv 1$, $b_n = 3 n^2 + n + 1$); the same computation
with symbolic $(\alpha_*, \beta_*)$ proves the universality
identity. Substituting $\theta = z\partial_z$ in
\eqref{eq:Lf-cc} with $(\alpha_1, \alpha_0, \beta_2, \beta_1,
\beta_0) = (0, 1, 3, 1, 1)$ yields the order-$2$ ODE
\[
  L f := f - z \bigl( 3(\theta+1)^2 + (\theta+1) + 1 \bigr) f
        - z^2 \, f \;=\; 1.
\]
The Newton polygon of the homogeneous operator is read off
from the lattice points
$\{(0, 0), (1, 0), (1, 1), (1, 2), (2, 0)\}$
(weight = order of $\theta$, height = order of $z$).
The lower convex hull has a single non-trivial edge
$(0, 0) \to (1, 2)$ of slope $1/2$, corresponding to a
Gevrey-$2$-in-$z$ irregular singularity of multiplicity $2$.
Substituting $z = u^2$, $\theta = (u/2)\partial_u$, the
operator becomes Gevrey-$1$-in-$u$, and the level-$1/u$
trans-series ansatz $f(u) = \exp(c/u)\,u^{\rho}\,\sum_{k \ge 0} a_k u^k$
produces the characteristic polynomial
$\chi(c) = 1 - (\beta_2/4) c^2 = 1 - 3 c^2 / 4$ at leading
order. The two roots $c = \pm 2/\sqrt{3}$ are the
Borel-plane singular distances of the formal pair $f_\pm$.
The $u^1$ coefficient of the trans-series equation pins
$\rho$ via the indicial polynomial; for general
$(\beta_2, \beta_1)$ this gives
$\rho = -3/2 - \beta_1/\beta_2$, and for V$_\mathrm{quad}$
specifically $\rho = -11/6$. The Birkhoff recursion for the
$a_k$ then proceeds row-by-row in the $u^k$ coefficient and
is carried to $K = 200$ in CC-PIPELINE-G
(\texttt{newton\_birkhoff.py}, \texttt{vquad\_recovery.py}).

The two key observations from the worked example are
structural:

\begin{enumerate}
\item The $\alpha_1, \alpha_0$ entries enter the operator
only at the $z^2$ Newton-polygon vertex $(2, 0)$, i.e.\ at
order $u^4$ in the $u$-uniformised operator. Therefore
$\xi_0$ (read off at order $u^0$) and $\rho$ (read off at
order $u^1$) are independent of the numerator. This is the
structural reason $\xi_0$ is $\beta_2$-only.
\item The first $a_k$ that depends on $\alpha_1$ is $a_2$,
entering at order $u^2$ in the trans-series equation, where
the $z^2 = u^4$ piece of the operator first reaches the
running coefficient. This is the structural reason
$\Cref{prop:noninj}$ separates the QL01/QL02 pair at $a_2$
and not before.
\end{enumerate}

Numerically, on V$_\mathrm{quad}$:
$a_1 \approx 0.6374909222302118$,
$a_2 = 0.4792390046296\overline{296}$ (exact rational
$\frac{2585}{5394}$ at low order, full exact form in
\cite{siarc_cc_pipeline_g}),
$a_3 \approx 0.7093200805668465$,
$a_4 \approx 1.1264939807461013$, agreeing with the
literature $\Vquad$ formal solution to all $200$ digits of
CC-PIPELINE-G's working precision.

\paragraph{Why the BoT space failed and the $\CC$ space succeeded.}
The UF-1 finding of CC-PIPELINE-F
\cite{siarc_cc_pipeline_f} is now structurally explained:
the BoT channel takes its formal series in $1/n$ from the
trans-series of $\log|\dn|$, whose Borel transform is the
image of a non-linear (logarithmic) gauge applied to the
same underlying analytic data. The $\CC$ channel takes its
formal series in $u = \sqrt z$ directly from the
generating-function ODE \eqref{eq:Lf-cc}; the Newton polygon
at $z = 0$ is read off as a polytope of the operator and
carries the irregular singular structure linearly. The two
spaces are related by an asymptotic transform (essentially a
Mellin--Borel composite) that is non-trivial in general; the
empirical $-0.5$-digit failure of
\cite{siarc_cc_pipeline_f} is the failure of that transform
to preserve the Borel singular distance. This is consistent
with the position taken in \Cref{rem:channel-vs-sum} that
two channels can use closely related summation methods and
still carry different singular data.

\Cref{thm:vquad-cc} closes the keystone test of v1.0:
the $\CC$ channel does recover the $\Vquad$ reduction, and
does so structurally rather than numerically. The remaining
Laplace gap is the refined open problem
\textsf{op:cc-formal-borel} (\Cref{sec:open}).

\subsection{Comparison table}
\label{ssec:comparison}

\begin{center}
\renewcommand{\arraystretch}{1.25}
\begin{tabular}{lcccc}
\toprule
Channel & Space $D$ & Gauge $T$ & Section $S$ &
  Detects $\Vquad$? \\
\midrule
$\Lt$  & series in $1/n$, $t$-deformed &
  uniform on $t$-grid &
  Borel--Laplace & no \\
$\BoT$ & $1/n$ trans-series of $\log|\dn|$ &
  Pad\'e pole-radius stability &
  Pad\'e--Borel & no \\
$\CC$  & formal sol.\ of \eqref{eq:Lf-cc} in $u=\sqrt z$ &
  $(\xi_0,\rho,a_k)$ exactness &
  Riemann--Hilbert & yes (\Cref{thm:vquad-cc}) \\
\bottomrule
\end{tabular}
\end{center}

\subsection{The locked WKB exponent identity}
\label{ssec:wkb}

\begin{proposition}[WKB exponent identity, empirical;
\cite{siarc_pcf1_v13,siarc_borel_kscan,siarc_cc_pipeline_f}]
\label{prop:wkb}
For each of the six $\Delta<0$ degree-$2$ families surveyed,
the leading trans-series exponents of
$\log|\dn| = -A n \log n + \alpha n - \beta \log n + \gamma
+ \sum_{k \ge 1} h_k / n^k$
satisfy the closed form
\[
  A \;=\; \begin{cases} 4 & a_n \equiv \text{const}, \\
                        3 & a_n = c_a n + d_a, \end{cases}
  \qquad
  \alpha \;=\; A \;-\; 2 \log c_b \;+\; \log|c_a|,
\]
verified to $\geq 13$ decimal digits in
\cite{siarc_borel_kscan} and to $\ge 23$ digits in
five families in \cite{siarc_cc_pipeline_f}.
\end{proposition}

\Cref{prop:wkb} fixes the leading WKB exponent of
$\log|\dn|$ but \emph{not} the residual Laplace step that
converts the Birkhoff-factored formal pair $f_\pm(u)$ of the
$\CC$ channel into Stokes data; that step is formally
separate, and Theory-Fleet H4 (see \Cref{rem:h4-median}
of \S\ref{ssec:h4-median}, and the new
\textsf{op:cc-median-resurgence-execute} of
\Cref{sec:open}) predicts that median \'Ecalle resurgence
on the V$_{\mathrm{quad}}$ formal coefficient series at
$5000$ coefficients / dps $250$ extracts the alien
amplitude at $\zeta_*=4/\sqrt 3$ to $30$--$50$ digits.

The relevance of \Cref{prop:wkb} for v1.1 is the
cross-reference to the cubic case
$A = 2d = 6$ \cite{siarc_pcf2_v11}: the degree-$2$ split
$A \in \{3, 4\}$ may be a low-degree artefact and the
universal form $A = 2d$ may hold for all $d \ge 3$ (PCF-2
\textsf{op:d2-anomaly}). At v1.2 the quartic confirmation
$A = 2d = 8$ ($60/60$ at $d=4$, PCF2-SESSION-Q1
\cite{siarc_pcf2_session_q1}) supports this reading. This in
turn motivates the $\xi_0$-extension problem
\textsf{op:xi0-degree-d} (\Cref{sec:open}), now empirically
resolved at $d=2,4$ via the linear formula
$\xi_0 = d / \beta_d^{1/d}$ (\Cref{conj:xi0-univ}).

% =====================================================================
\subsection{Median \'Ecalle resurgence as the residual
Laplace step (v1.3)}
\label{ssec:h4-median}

The $\CC$-channel recovery of $\Vquad$
(\Cref{thm:vquad-cc}) is exact-modulo-Laplace at the
algebraic level: $\xi_0 = 2/\sqrt 3$ and $\rho = -11/6$ are
verified to $200$ digits, but the connection-coefficient
ratio still requires a numerical Borel--Laplace summation
of the formal coefficient series $\sum_{k} a_k u^k$ to
extract the Stokes constant attached to the Borel
singularity at $\zeta_* = 4/\sqrt 3$. The Pad\'e--Borel
estimator used in CC-PIPELINE-G \cite{siarc_cc_pipeline_g}
showed slow convergence (Domb--Sykes secondary metric drift
in the leading digits), consistent with global
approximation to a branch cut rather than lack of
information in the coefficients. Theory-Fleet H4 of
2026-05-01 \cite{siarc_theory_fleet_h4} identifies the
correct local object as the alien derivative
$\Delta_{\zeta_*}$ in the sense of \'Ecalle--Sauzin
\cite{ecalle1981fonctions, sauzin2014resurgent,
mitschi2016divergent} and the correct unambiguous real
resummation as the median sum
$\mathcal S_{\mathrm{med}}$.

\begin{remark}[Median resurgence prediction; theoretical, H4]
\label{rem:h4-median}
Theory-Fleet H4 \cite{siarc_theory_fleet_h4} is a
\emph{theoretical prediction at HIGH confidence} -- not an
executed measurement -- that median \'Ecalle resurgence on
the V$_{\mathrm{quad}}$ formal coefficient series, run at
$5000$ coefficients with mpmath working precision
$\mathrm{dps} = 250$ and $\zeta_* = 4/\sqrt 3$ pinned to
$200$ digits, will extract the leading alien amplitude
$S_{\zeta_*}$ at $\zeta_*$ to $30$--$50$ digits, with a
central forecast of $\sim 40$ digits. The recipe is: fit
the branch exponent from the large-order tail by
Borel--Darboux relations, extract the Stokes coefficient
from the half-Stokes prescription
$\mathcal S_{\mathrm{med}} = \mathcal S_{0^-} \circ
\mathfrak S_0^{1/2}$, and cross-check via local Borel
singular-germ fitting in a neighbourhood of $\zeta_*$.
H4 explicitly does \emph{not} forecast a closed-form
$S_{\zeta_*}$ (Gamma products at thirds/sixths are
plausible candidates but the literature alone does not
prove the form). The sub-problem of executing the
prediction is logged as
\textsf{op:cc-median-resurgence-execute}
(\Cref{sec:open}, new at v1.3); resolving it would close
the residual Laplace step of \textsf{op:cc-formal-borel}
for $V_{\mathrm{quad}}$ and would conditionally close the
$\Vquad \rightsquigarrow P_{III}(D_6)$ correspondence at
the level of Stokes data, modulo a documented
normalisation map to the $P_{III}(D_6)$ isomonodromic
Stokes datum of \cite{lodayrichaud2016divergent,
boalch2001stokes}.
\end{remark}

% =====================================================================
\subsection{The cross-degree invariant triple of the SIARC
umbrella (v1.3)}
\label{ssec:invariant-triple}

The post-T2 refactor of the SIARC umbrella v2.0
\cite{siarc_umbrella_v2}, \S4.4 packages the cross-degree
SIARC stack as a \emph{bridge functor}
\[
  \Phi : \PCF(1,b) \;\longmapsto\;
  \bigl(\Delta_d(b),\;\|\Delta\|_{\mathrm{Pet}}(\tau_b),\;
        \xi_0(b)\bigr),
\]
the \emph{invariant triple}: the polynomial discriminant
$\Delta_d(b) = \mathrm{disc}(b) \in \mathbf{Q}^\times$, the
Petersson-normalised modular discriminant
$\|\Delta\|_{\mathrm{Pet}}(\tau_b) =
|\Delta(\tau_b)| \cdot (\mathrm{Im}\,\tau_b)^6$ at the
period $\tau_b$ of the curve $E_b: y^2 = b(x)$ (cubic case
$d=3$), and the Borel radius $\xi_0(b) = d/\beta_d^{1/d}$
of \Cref{conj:xi0-univ}. The Borel-radius axis that this
outline has carried since v1.0 is, in this language, the
third coordinate of $\Phi$.

The empirical content of the modular-discriminant axis is
the headline of PCF-2 v1.3 \cite{siarc_pcf2_v13}, T2 phase
B/C: on the deep R1.3 cubic catalogue ($n=50$ irreducible
non-singular cubic families, deep WKB at
$\mathrm{dps} = 4000$, $N_{\max} = 480$), the Petersson
height $\log\|\Delta\|_{\mathrm{Pet}}(\tau_b)$ correlates
with the residual $\delta = A_{\mathrm{fit}} - 6$ at
Spearman
\[
  \rho\bigl(\log\|\Delta\|_{\mathrm{Pet}}(\tau_b),\,\delta_{\mathrm{deep}}\bigr)
   \;=\; +0.638,
  \qquad p_{\mathrm{Bonf}}(K{=}14) = 8.6\times 10^{-6},
\]
\emph{beating} the bare-$\log|j|$ baseline
($\rho = -0.568$, $p_{\mathrm{Bonf}} = 2.34\times 10^{-4}$)
by a factor ${\sim}30$ in Bonferroni $p$
\cite{siarc_pcf2_session_t2}. Bare $\log|E_4(\tau_b)|$
alone underperforms $\log|j|$, consistent with the
modular identity $1728\,\Delta = E_4^3 - E_6^2$: at the
equianharmonic ($j = 0$) cell $E_4$ has a simple zero so
$\log|E_4|$ is locally pinned, while
$\log\|\Delta\|_{\mathrm{Pet}}
= 6\log\mathrm{Im}\,\tau + \log|\Delta|$ remains
well-conditioned globally.

Channel Theory adds nothing structurally new on either the
$\Delta_d$ or the $\|\Delta\|_{\mathrm{Pet}}$ axis: those
are PCF-1 / PCF-2 / umbrella-v2.0 territory
\cite{siarc_pcf1_v13, siarc_pcf2_v13, siarc_umbrella_v2}.
The contribution of the present outline is to identify the
$\xi_0$ axis component of $\Phi$ -- proven at $d=2$
(\Cref{prop:xi0}), verified at $d=4$ (\Cref{prop:xi0-d4}),
conjectural at $d \ne 2,4$ (\Cref{conj:xi0-univ}) -- and
to operationalise it through the channel functor $\chi$ of
\S\ref{sec:functor}. As a bookkeeping note, the bridge
conjecture's three observed tiers B1/B2/B3
(\Cref{conj:bridge}) are now naturally re-tiered against
the full triple: a
$(\beta_d,\,\|\Delta\|_{\mathrm{Pet}},\,\xi_0)$-coordinatised
refinement is the structurally correct framing, and the
$d=3$ Petersson-height signal of T2 is a missing
ingredient that the umbrella v2.0 now supplies. We do
\emph{not} restate B1/B2/B3 as theorems here; we simply
record the compatibility expectation that any successful
bridge identity $\Phi_b$ will factor through $\Phi$ on the
relevant axes.

% =====================================================================
\section{The channel functor}
\label{sec:functor}

The Master Conjecture \cite{siarc_umbrella_v2} posits a functor
$\chi : \Sigmaspace \to \Channel$ where $\Sigmaspace$ is a
category of deformations of PCF families. We make
$\Sigmaspace$ slightly more precise.

\begin{definition}[Object of $\Sigmaspace$]
An object of $\Sigmaspace$ is a triple
$\Sigma = (F, \omega, h)$ where $F = \PCF(a, b)$ is a PCF
family as in the scope of \Cref{sec:position},
$\omega \in \mathbf{Z}[n]$ is a deformation direction, and
$h > 0$ is a step parameter. The deformed family is
$F_t = \PCF(a, b + t \omega)$ for $|t| < h$.
\end{definition}

\begin{definition}[Morphism of $\Sigmaspace$]
A morphism
$(F, \omega, h) \to (F', \omega', h')$ exists when $F = F'$,
$\omega' = \lambda \omega$ for $\lambda \in \mathbf{Q}^\times$,
and $h' \le h / |\lambda|$. Composition is multiplication of
the scalars $\lambda$.
\end{definition}

\begin{conjecture}[Properties of $\chi$]
\label{conj:chi}
The channel functor satisfies:
\begin{enumerate}[label=(\roman*)]
\item Functoriality: changing the deformation $\Sigma$ on a
fixed family genuinely changes the target channel $\chi(\Sigma)$.
\item Empirical consistency:
$\chi(\Sigma(\Vquad)) = \CC$ (now a theorem-modulo-Laplace,
\Cref{thm:vquad-cc}).
\item Non-constancy: there exist $\Sigma_1 \neq \Sigma_2$
with $\chi(\Sigma_1) \ne \chi(\Sigma_2)$.
\end{enumerate}
\end{conjecture}

\Cref{conj:chi}(ii) is the substantive change at v1.1: where
v1.0 had a single empirical data point ($\Vquad$ in $\CC$),
v1.1 has a structural identity (\Cref{prop:xi0}) plus an
exact algebraic recovery (\Cref{thm:vquad-cc}) on which
(ii) rests.

% =====================================================================
\section{Conjectured no-go theorem}
\label{sec:nogo}

\begin{conjecture}[$\Lt$ no-go for degree $2$, $\Delta<0$]
\label{conj:nogo}
Let $F = \PCF(a, b)$ be a polynomial continued fraction with
$a, b \in \mathbf{Z}[n]$, $b$ irreducible and
$\mathbf{Z}$-primitive, $\deg b = 2$, and discriminant
$\Delta(b) < 0$. Then no Painlev\'e reduction of the
$\Lt$-deformation $F_t = \PCF(a, b + t\omega)$ exists in the
$\Lt$ channel of \Cref{ex:lt}.
\end{conjecture}

\paragraph{Empirical evidence and falsifiability.}
\Cref{conj:nogo} is consistent with the Session~D record
\cite{siarc_painleve_deep_6x}: 0/6 families admit a Painlev\'e
fit at residual $\le 10^{-4}$, with residuals stable under
precision escalation. It is falsified by any single
$\Delta < 0$ degree-$2$ family for which an $\Lt$-channel
Painlev\'e fit residual drops below $10^{-4}$ stably.
\Cref{prop:xi0} does not bear directly on \Cref{conj:nogo} --
the two channels have different formal-series spaces -- but
\Cref{rem:kscan} sharpens the failure mode: same $\beta_2$
forces same $\xi_0$ in $\CC$ without forcing same $\Lt$
Painlev\'e label.

% =====================================================================
\section{Bridge conjecture (revised, v1.1)}
\label{sec:bridge}

The v1.0 outline parametrised the bridge identity by the
discriminant $\Delta$. \Cref{prop:xi0,prop:noninj} show
that the discriminant is the wrong invariant: the structure
that $\CC$ detects depends on the full coefficient tuple, with
$\beta_2$ alone fixing $\xi_0$ but not the Painlev\'e label,
and the secondary tower $(\rho, a_1, a_2, \ldots)$ adding
discrimination.

\begin{conjecture}[Bridge conjecture, revised]
\label{conj:bridge}
For every degree-$2$ PCF $(1, b)$ in scope, there exists a
\emph{bridge identity} $\Phi_b$ mapping the $\CC$-channel
datum $(\xi_0, \rho, a_1, a_2, \ldots)$ to a Painlev\'e
$P$-class label, with $\Phi_b$ functorial in the coefficient
tuple $(\alpha_1, \alpha_0, \beta_2, \beta_1, \beta_0)$
rather than in any single invariant.
The bridge identity comes in three observed variants:
\begin{enumerate}[label=B\arabic*.]
\item \emph{Single-coefficient witness.} Painlev\'e label
determined by $\beta_2$ alone, witnessed by
$\xi_0 = 2/\sqrt{\beta_2}$ (\Cref{prop:xi0}). This tier
contains $\Vquad$ in the trivial sense that the
\emph{singularity radius} alone is $\beta_2$-determined.
\item \emph{Two-coefficient witness.} Painlev\'e label
determined by $(\beta_2, \text{secondary invariant})$;
the secondary invariant is enough to distinguish QL01 from
QL02 (which share $\xi_0$ and $\rho$ and split at $a_2$,
\Cref{prop:noninj}).
\item \emph{Full-tuple witness.} Painlev\'e label requires
the entire coefficient tuple
$(\alpha_1, \alpha_0, \beta_2, \beta_1, \beta_0)$.
\end{enumerate}
\end{conjecture}

\paragraph{The B1 case for $\Vquad$.} For $\Vquad$,
$\beta_2 = 3$, so by \Cref{prop:xi0} $\xi_0 = 2/\sqrt 3$,
and by \Cref{thm:vquad-cc} the $P_{III}(D_6)$ label is
recovered. Whether $\xi_0$ alone suffices to predict
$P_{III}(D_6)$ for an unseen family with $\beta_2 = 3$ is
an open question: $Q_{L15}$ also has $\beta_2 = 3$ but is a
K-SCAN artefact at the BoT level (\Cref{rem:kscan}) and the
fingerprint agreement at $(\rho, a_k)$ is $0.02$ digits, so
it does \emph{not} share $\Vquad$'s Painlev\'e label.

\paragraph{Tier discrimination is open.} The K-SCAN data does
not yet discriminate B1 vs.\ B2 vs.\ B3 across the full
catalogue: even where the Painlev\'e label is unknown
(QL01--QL26 minus $\Vquad$), the bridge tier remains
unspecified. This is the new
\textsf{op:bridge-witness-tier} of \Cref{sec:open}.

\paragraph{Per-family tier candidacy.}
\Cref{tab:bridge-tier} records, for each of the six
$\Delta < 0$ degree-$2$ families, which bridge tier the
family is currently \emph{compatible with}, given
\Cref{tab:cc-channel-v2}. A tier-B1 candidate is a family
that shares $\beta_2$ with another family of known Painlev\'e
label (only $\Vquad$ has a known label, so this is
``$\beta_2 = 3$''). A tier-B2 candidate is a family that
shares both $\beta_2$ and $\beta_1$ (equivalently $\xi_0$ and
$\rho$) with another family of known label. A tier-B3
candidate is any family not in B1 or B2's candidate set.

\begin{table}[ht]
\centering
\small
\caption{Per-family bridge-tier candidacy under
\Cref{conj:bridge}. ``Tier'' is the lowest tier compatible
with the present data; ``predicted label'' is what tier B1
\emph{would} predict if $\beta_2$ alone determined the
Painlev\'e label, and ``actual label''
is what is currently known. The mismatch on $Q_{L15}$ is
\Cref{rem:kscan}.}
\label{tab:bridge-tier}
\begin{tabular}{l c c c c c}
\toprule
Family & $\beta_2$ & $\rho$ & Tier candidate &
  Predicted (B1) & Known label \\
\midrule
$\Vquad$  & $3$ & $-11/6$ & B1 anchor   & $P_{III}(D_6)$ &
   $P_{III}(D_6)$ \\
$Q_{L01}$ & $1$ & $-5/2$  & B2 candidate & --              & unknown \\
$Q_{L02}$ & $1$ & $-5/2$  & B2 candidate & --              & unknown \\
$Q_{L06}$ & $2$ & $-1$    & B3 candidate & --              & unknown \\
$Q_{L15}$ & $3$ & $-5/6$  & B1 reject    & $P_{III}(D_6)$ &
   $\ne P_{III}(D_6)$ \\
$Q_{L26}$ & $4$ & $-1$    & B3 candidate & --              & unknown \\
\bottomrule
\end{tabular}
\end{table}

The $Q_{L15}$ row is the only family in the roster where B1
is already \emph{ruled out}: shared $\beta_2 = 3$ with
$\Vquad$ predicts $P_{III}(D_6)$ at the B1 level, but the
$(\rho, a_k)$-tower agreement is $0.02$ digits
(\Cref{tab:cc-channel-v2}), so the prediction fails. This
pushes the $\Vquad$ witness for B1 down to a B2 or higher
statement, depending on \textsf{op:cc-monodromy-RH}. The
QL01/QL02 pair is the cleanest B2 candidate: identical
$(\xi_0, \rho)$, splitting only at $a_2$
(\Cref{prop:noninj}); whether they share a Painlev\'e label
or not is exactly the B2-vs-B3 question.

\paragraph{Three structural variants of the existence statement.}
We retain the v1.0 distinction between
\emph{existence} (some $\Phi_b$ exists),
\emph{algebraicity} ($\Phi_b$ algebraic in its arguments),
and \emph{approximate} (a tolerance $\varepsilon$ relaxation).
Existence is the minimum, algebraicity the aspiration,
approximate the fallback.

\paragraph{Obstructions remaining.}
\Cref{prop:xi0,prop:noninj,thm:vquad-cc} dispose of the v1.0
obstruction that ``the Borel singular points $\zeta_1, \ldots,
\zeta_r$ are not well-defined for the five non-$\Vquad$
families''. They are now well-defined, exact, and tabulated
(\Cref{tab:cc-channel-v2}). The remaining obstruction is the
absence of a Painlev\'e label for the five non-$\Vquad$
families: $\Phi_b$ cannot be tested without a target.

% =====================================================================
\section{Related work and prior channel structures}
\label{sec:related}

\subsection{Resurgence and alien calculus}
The most developed pre-existing analogue is the
resurgence-theoretic framework of
\cite{ecalle1981fonctions, costin2008asymptotics,
sauzin2014resurgent, mitschi2016divergent}. Our $\BoT$
channel is essentially a discrete Pad\'e-truncated version;
the failure mode of \cite{siarc_borel_5x} is consistent with
Pad\'e approximation not always recovering the alien
structure. The continuous-side success of resurgence
\cite{delabaere2016divergent} on $P_I$ remains the structural
benchmark.

\subsection{Stokes / isomonodromic theory}
The $\CC$ channel has a long history in the theory of
isomonodromic deformations \cite{jimbomiwa1981monodromy,
vanderputsinger2003galois, boalch2001stokes}. The
operationalisation via Newton polygon plus Birkhoff
\cite{siarc_cc_pipeline_g} is the discrete-difference
analogue.

\subsection{Heun--Painlev\'e correspondence}
\cite{slavyanov2000special} treats Heun-class linear ODEs and
their reductions to Painlev\'e. The $\Vquad$ Casoratian
specialises in this framework, and that specialisation is
how the $(1/6, 0, 0, -1/2)$ moduli point is computed.

\subsection{Algebraic Painlev\'e moduli}
\cite{mazzocco2001picard, mazzocco2002painleve, iwasaki1991from,
ohyama2006studies, guzzetticlass2012} give finite catalogues
of algebraic and Picard-type solutions. The T2B /
UMB-PVI-MATCH umbrella result \cite{siarc_t2b_v3} reads PCF
families into this catalogue: Class B 16/16 lands in
$P_{III}(D_6)$, with one $P_{VI}$-Picard rigidity candidate
flagged.

\subsection{SIARC sessions cited at v1.1}
The v1.1 outline is empirically anchored in:
\begin{itemize}
\item PCF-1 v1.3 \cite{siarc_pcf1_v13} (concept DOI
\href{https://doi.org/10.5281/zenodo.19937196}{10.5281/zenodo.19937196})
-- WKB exponent identity, $\Lt$-channel Stokes signal,
$\Vquad / P_{III}(D_6)$ reduction.
\item PCF-2 v1.1 \cite{siarc_pcf2_v11} (concept DOI
\href{https://doi.org/10.5281/zenodo.19939463}{10.5281/zenodo.19939463})
-- Conjecture~B4 ($A = 2d$), \textsf{op:d2-anomaly},
cross-reference for \textsf{op:xi0-degree-d}.
\item PCF2-SESSION-Q1 \cite{siarc_pcf2_session_q1}
(SIARC bridge,
\href{https://github.com/papanokechi/siarc-relay-bridge/tree/main/sessions/2026-05-01/PCF2-SESSION-Q1/}{sessions/2026-05-01/PCF2-SESSION-Q1/})
-- quartic harvest at $d=4$: $A = 2d = 8$ verified $60/60$
and the universal Newton-polygon characteristic root
$\xi_0(b) = d / \beta_d^{1/d}$ verified at $d=4$ across $8$
representatives with spread $0$
(\Cref{prop:xi0-d4,conj:xi0-univ}).
\item BOREL-CHANNEL-K-SCAN \cite{siarc_borel_kscan} -- BoT
channel artefact diagnosis (bridge URL
\href{https://github.com/papanokechi/siarc-relay-bridge/tree/main/sessions/2026-05-01/BOREL-CHANNEL-K-SCAN/}{sessions/2026-05-01/BOREL-CHANNEL-K-SCAN/}).
\item CC-PIPELINE-F \cite{siarc_cc_pipeline_f} -- UF-1
formal-series-space disagreement (bridge URL
\href{https://github.com/papanokechi/siarc-relay-bridge/tree/main/sessions/2026-05-01/CC-PIPELINE-F/}{sessions/2026-05-01/CC-PIPELINE-F/}).
\item CC-PIPELINE-G \cite{siarc_cc_pipeline_g} --
Newton-polygon $+$ Birkhoff recovery, three universality
findings (bridge URL
\href{https://github.com/papanokechi/siarc-relay-bridge/tree/main/sessions/2026-05-01/CC-PIPELINE-G/}{sessions/2026-05-01/CC-PIPELINE-G/}).
\end{itemize}

The novel structural content of v1.1 is
\Cref{prop:xi0,prop:noninj,thm:vquad-cc} and the
B1/B2/B3 tier reformulation of \Cref{conj:bridge}.

\subsection{Additional sessions cited at v1.3}
\label{ssec:sessions-v13}
v1.3 additionally cites the post-T2 cohort of 2026-05-01 /
2026-05-02:
\begin{itemize}
\item SIARC umbrella v2.0 \cite{siarc_umbrella_v2} (concept
DOI \href{https://doi.org/10.5281/zenodo.19885549}{10.5281/zenodo.19885549};
v2.0 at \href{https://doi.org/10.5281/zenodo.19965041}{10.5281/zenodo.19965041})
-- post-T2 cross-degree refactor of the umbrella around
the invariant triple
$(\Delta_d,\,\|\Delta\|_{\mathrm{Pet}},\,\xi_0)$;
\S4.4 supplies the cite handle for the $\xi_0$ axis.
\item PCF-2 v1.3 \cite{siarc_pcf2_v13} (concept DOI
\href{https://doi.org/10.5281/zenodo.19936297}{10.5281/zenodo.19936297};
v1.3 at \href{https://doi.org/10.5281/zenodo.19963298}{10.5281/zenodo.19963298})
-- cubic-modular framing absorbed at v1.3; T2 Petersson
height result at $\rho = +0.638$ on $n=50$ cubic families.
\item PCF2-SESSION-T2 \cite{siarc_pcf2_session_t2}
(bridge URL
\href{https://github.com/papanokechi/siarc-relay-bridge/tree/main/sessions/2026-05-02/PCF2-SESSION-T2/}{sessions/2026-05-02/PCF2-SESSION-T2/})
-- Petersson-height correlation at
$\rho = +0.638$, $p_{\mathrm{Bonf}} = 8.6\times 10^{-6}$
on $n=50$ cubic families.
\item THEORY-FLEET H4 \cite{siarc_theory_fleet_h4}
(bridge URL
\href{https://github.com/papanokechi/siarc-relay-bridge/tree/main/sessions/2026-05-01/THEORY-FLEET/H4/}{sessions/2026-05-01/THEORY-FLEET/H4/})
-- HIGH-confidence theoretical prediction
\textsf{MEDIAN\_RESURGENCE\_GIVES\_30+\_DIGITS} for
$V_{\mathrm{quad}}$ at $\zeta_* = 4/\sqrt 3$.
\item THEORY-FLEET AGGREGATOR (bridge URL
\href{https://github.com/papanokechi/siarc-relay-bridge/tree/main/sessions/2026-05-01/THEORY-FLEET/AGGREGATOR/}{sessions/2026-05-01/THEORY-FLEET/AGGREGATOR/})
-- unified verdict aggregating H1--H6 of the multi-agent
theoretical fleet of 2026-05-01.
\item PCF2-V13-RELEASE (bridge URL
\href{https://github.com/papanokechi/siarc-relay-bridge/tree/main/sessions/2026-05-02/PCF2-V13-RELEASE/}{sessions/2026-05-02/PCF2-V13-RELEASE/})
and SIARC-UMBRELLA-V2-RELEASE (bridge URL
\href{https://github.com/papanokechi/siarc-relay-bridge/tree/main/sessions/2026-05-02/SIARC-UMBRELLA-V2-RELEASE/}{sessions/2026-05-02/SIARC-UMBRELLA-V2-RELEASE/})
-- companion release sessions on 2026-05-02 producing
PCF-2 v1.3 and SIARC umbrella v2.0 respectively, whose
content this v1.3 outline absorbs.
\end{itemize}

% =====================================================================
\section{Implications for the Master Conjecture}
\label{sec:implications}

A rigorous announcement of the Master Conjecture requires:
\begin{enumerate}
\item Pipeline closure for the five non-$\Vquad$ families
(\textsf{op:cc-monodromy-RH}). \Cref{prop:xi0,tab:cc-channel-v2}
provide the formal-solution data for all six families;
the missing piece is the full Stokes-matrix datum.
\item No-go proof or refutation (\textsf{op:no-go-proof})
of \Cref{conj:nogo}.
\item Bridge identity tier discrimination
(\textsf{op:bridge-witness-tier}) for at least one family
beyond $\Vquad$.
\item Cubic extension along
\cite{siarc_pcf2_v11, siarc_pcf2_program}
(\textsf{op:xi0-degree-d}).
\end{enumerate}

% =====================================================================
\section{Open problems and program}
\label{sec:open}

The v1.0 list is restructured: the previous keystone
\textsf{op:cc-pipeline-keystone} is \emph{resolved} by
\Cref{thm:vquad-cc}; the residual Laplace step is split out
as \textsf{op:cc-formal-borel}; new problems
(\textsf{op:xi0-degree-d}, \textsf{op:xi0-d3-direct},
\textsf{op:bridge-witness-tier},
\textsf{op:coefficient-tuple-functoriality}) are added.
Version~1.2 narrows \textsf{op:xi0-degree-d} to its
direct-$d=3$ component (\textsf{op:xi0-d3-direct}) after the
PCF2-SESSION-Q1 confirmation at $d=4$; the cross-referenced
PCF-2 problem \textsf{op:b4-degree-d} \cite{siarc_pcf2_v11}
is correspondingly narrowed to $d \ge 5$ (PCF-2's
$A = 2d = 8$ is verified $60/60$ at $d=4$ in
PCF2-SESSION-Q1 \cite{siarc_pcf2_session_q1}).

\begin{problem}[\textsf{op:cc-median-resurgence-execute} (new at v1.3)]
Run median \'Ecalle resurgence on the V$_{\mathrm{quad}}$
formal coefficient series at $5000$ coefficients with
mpmath working precision $\mathrm{dps} = 250$, with
$\zeta_* = 4/\sqrt 3$ pinned to $200$ digits. Extract the
leading alien amplitude $S_{\zeta_*}$ at $\zeta_*$, fit the
branch exponent from the large-order tail by Borel--Darboux
relations, and cross-check via local Borel singular-germ
fitting in a neighbourhood of $\zeta_*$. Theory-Fleet H4
\cite{siarc_theory_fleet_h4} \emph{predicts at HIGH
confidence} a $30$--$50$ digit precision (central forecast
${\sim}40$ digits) for $S_{\zeta_*}$. Resolves the
residual Laplace step of \textsf{op:cc-formal-borel} for
V$_{\mathrm{quad}}$ and would conditionally close the
$\Vquad \rightsquigarrow P_{III}(D_6)$ correspondence at
the Stokes-data level under a documented normalisation map.
\emph{Tractability.} $1$--$3$ sessions of compute (mpmath
$+$ a custom alien-derivative extractor); the bottleneck is
the $5000$-coefficient generation, not the alien-amplitude
extraction.
\end{problem}

\begin{problem}[\textsf{op:cc-formal-borel}]
\textbf{Status (v1.3):} \emph{PARTIALLY DIAGNOSED} by
Theory-Fleet H4 \cite{siarc_theory_fleet_h4}, which
predicts at HIGH confidence that median \'Ecalle
resurgence at $5000$ coefficients / $\mathrm{dps}\,250$
will extract the alien amplitude at $\zeta_* = 4/\sqrt 3$
to $30$--$50$ digits. The residual question is now the
\emph{execution} of that prediction
(\textsf{op:cc-median-resurgence-execute}, new at v1.3),
not its formulation.

Numerically Borel--Laplace sum the V$_\mathrm{quad}$
formal coefficient series $\sum_k a_k u^k$ at convergent
tail accuracy. The Domb--Sykes secondary metric in
CC-PIPELINE-G \cite{siarc_cc_pipeline_g} showed slow
convergence; produce a clean numerical proof-of-concept of
the connection-coefficient ratio at $\ge 30$ digits. The
algebraic identities $\xi_0 = 2/\sqrt 3$, $\rho = -11/6$
are already at $200$ digits; the Laplace step is the residual.
\emph{Tractability.} 1--2 sessions for a Borel--\'Ecalle
ray-summation test on $\Vquad$.
\end{problem}

\begin{problem}[\textsf{op:xi0-degree-d}]
Extend \Cref{prop:xi0} to PCFs of arbitrary degree $d$. For
degree-$d$ denominator $b(n) = \beta_d n^d +
\mathrm{l.o.t.}$, conjecture $\xi_0 = c(d) / \beta_d^{1/d}$
(\Cref{conj:xi0-univ}): candidate $c(d) = d$ (linear),
proven at $d=2$ (\Cref{prop:xi0}), verified at $d=4$
with spread $0$ in PCF2-SESSION-Q1
(\Cref{prop:xi0-d4}, \cite{siarc_pcf2_session_q1});
$d=3$ verification deferred to \textsf{op:xi0-d3-direct}.
The earlier v1.1 candidate $c(d) = 2\sqrt{(d-1)!}$ is
empirically falsified at $d=4$
(\Cref{rem:c-d-falsified}).
\emph{Tractability.} Resolved at $d=2,4$; remaining work is
the direct $d=3$ test (\textsf{op:xi0-d3-direct}) and a
general-$d$ Newton-polygon proof.
\end{problem}

\begin{problem}[\textsf{op:xi0-d3-direct}]
\textbf{Status (v1.3):} Carried forward unchanged.
Tractability re-checked at $1$--$2$ hours of computation,
no halt.

Carry out the direct Newton-polygon test of
\Cref{conj:xi0-univ} at $d=3$: select one cubic
representative per Galois bin of PCF-2 v1.1
\cite{siarc_pcf2_v11} ($+_{S_3}$ real, $+_{C_3}$ real,
$-_{S_3}$ CM), extract the characteristic polynomial of the
irregular singular point of $L_3 f = (1 - z B_3(\theta+1) -
z^2) f$ at $z=0$, measure $\xi_0$ at $\mathrm{dps} \ge 80$,
and compare with the prediction $c(3) = 3$, i.e.\
$\xi_0 = 3 / \beta_3^{1/3}$. Evidence at $d=4$ has
$c(4) = 4$ with spread $0$
(\Cref{prop:xi0-d4}, \cite{siarc_pcf2_session_q1}).
\emph{Tractability.} 1--2 hours of computation along the
PCF2-SESSION-Q1 \texttt{session\_q1\_newton.py} pattern,
adapted to the cubic operator.
\end{problem}

\begin{problem}[\textsf{op:bridge-witness-tier}]
For each degree-$2$ PCF in the catalogue, determine whether
the bridge identity is B1 ($\beta_2$ only), B2 ($\beta_2$
plus one secondary), or B3 (full tuple). Currently $\Vquad$
and $Q_{L15}$ share $\beta_2 = 3$ and would lie in B1's
candidate set if their Painlev\'e labels matched; they do
not (K-SCAN artefact, \Cref{rem:kscan}), pushing the
$\Vquad$ witness toward B2 or higher.
\emph{Tractability.} Hinges on \textsf{op:cc-monodromy-RH}.

\textbf{(v1.3):} Tier discrimination is now also informed
by the umbrella's $\|\Delta\|_{\mathrm{Pet}}$ axis
\cite{siarc_pcf2_v13}; for the cubic catalogue the
modular-discriminant signal is at $\rho = +0.638$ ($n=50$,
Bonferroni $p = 8.6\times 10^{-6}$), suggesting a
$(\beta_d,\,\|\Delta\|_{\mathrm{Pet}},\,\xi_0)$-coordinatised
refinement of the B1/B2/B3 tiering.
\end{problem}

\begin{problem}[\textsf{op:coefficient-tuple-functoriality}]
Make the functoriality statement of \Cref{conj:bridge}
precise: define the category of admissible PCF coefficient
tuples and morphisms (rescalings, polynomial
reparametrisations) and prove (or disprove) that the bridge
map $\Phi_b$ is a functor. Currently stated only
conjecturally.
\emph{Tractability.} Multi-year.
\end{problem}

\begin{problem}[\textsf{op:no-go-proof}]
Prove or refute \Cref{conj:nogo}.
\emph{Tractability.} 1--3 sessions for a refutation if a
moduli of perturbation directions $\omega$ produces a stable
$\Lt$ Painlev\'e fit; multi-year for a proof.
\end{problem}

\begin{problem}[\textsf{op:cc-monodromy-RH}]
Extend the Newton-polygon / Birkhoff formal-solution extractor
of \cite{siarc_cc_pipeline_g} to a full Riemann--Hilbert
monodromy datum at the irregular singular point $z=0$ via
Stokes-matrix computation. This is the prerequisite for a
rigorous Painlev\'e-class verdict on the five non-$\Vquad$
families at any finite-digit threshold.
\emph{Tractability.} 1--2 sessions per family.
\end{problem}

\begin{problem}[\textsf{op:channel-moduli}]
Is the channel space $\Channel$ a moduli space (with
positive-dimensional families of channels) or a discrete /
finite set?
\emph{Tractability.} Multi-year. Likely requires the
Stokes-matrix / Frobenius-manifold framework of
\cite{boalch2001stokes}.
\end{problem}

\begin{problem}[\textsf{op:borel-channel}]
Is the BoT channel of \Cref{ex:bot} ANOMALOUS in the formal
sense (the Borel-resummation section $S$ on the BoT
trans-series space does not exist for $\Delta<0$ degree-$2$
PCFs), or merely under-resolved at $K \le 24$?
The CC-PIPELINE-F UF-1 finding
\cite{siarc_cc_pipeline_f} -- that the BoT space and the
$\CC$ space disagree at the level of the Borel singular
structure -- argues for the former; the K-SCAN drift of
\cite{siarc_borel_kscan} is consistent with either reading.
PCF-1 v1.3 \cite{siarc_pcf1_v13} carries this as
``\emph{BOREL CHANNEL CONFIRMED ANOMALOUS}''.
\emph{Tractability.} Bound to \textsf{op:cc-formal-borel}.
\end{problem}

\begin{problem}[\textsf{op:cc-channel-existence} (legacy)]
For families with $\deg b \ge 3$ or $\Delta \ge 0$, does the
generating-function ODE \eqref{eq:Lf-cc} (or its
generalisation) still admit a Newton-polygon Birkhoff
formal pair, i.e., does the $\CC$ channel exist as a non-empty
triple $(D, T, S)$?
\emph{Tractability.} Joint with \textsf{op:xi0-degree-d}.

\textbf{(v1.3):} The cubic case is now bound to PCF-2 v1.3
\cite{siarc_pcf2_v13} (Conjecture B5/B6 in their $d=3$
restricted form), which gives a structural footprint to
test in the cubic CC channel.
\end{problem}

\begin{problem}[\textsf{op:painleve-classification}]
For families admitting a $\CC$ formal pair, classify the
attainable Painlev\'e labels. Empirically, only $P_{III}(D_6)$
appears in the present roster (one family).
\emph{Tractability.} Bound to \textsf{op:cc-monodromy-RH}.
\end{problem}

\begin{problem}[\textsf{op:non-piii} (legacy)]
Are there degree-$2$ $\Delta < 0$ PCFs whose $\CC$ reduction
is non-$P_{III}$? PCF-1 v1.3 Sec.~8 redirects this
to \textsf{op:borel-channel} on the grounds that QL06 / P-V
was ruled out by Session~D and that
the Borel channel diagnoses the candidate set is
artefact-dominated.
\emph{Tractability.} Bound to \textsf{op:cc-monodromy-RH}.
\end{problem}

% =====================================================================
\section*{Acknowledgements}

We thank the SIARC relay infrastructure
(\href{https://github.com/papanokechi/siarc-relay-bridge}%
{papanokechi/siarc-relay-bridge}) for storing and
checkpointing the empirical record cited throughout.

\paragraph{AI disclosure.} This outline was drafted with the
assistance of large language models in the SIARC
self-iterating analytic relay chain, including
Sessions~PAINLEVE-DEEP-6X, BOREL-CHANNEL-5X,
BOREL-CHANNEL-K-SCAN, CC-PIPELINE-F (UF-1) and CC-PIPELINE-G
(Newton-polygon $+$ Birkhoff recovery, three universality
findings). All numerical claims trace to exact computations
stored under \texttt{sessions/2026-05-01/} on the SIARC
bridge.

The v1.3 revision additionally absorbs the SIARC umbrella
v2.0 invariant-triple framing \cite{siarc_umbrella_v2}, the
PCF-2 v1.3 cubic-modular Petersson-height result
\cite{siarc_pcf2_v13}, and the Theory-Fleet H4
median-resurgence prediction
\cite{siarc_theory_fleet_h4}. AI tools (GitHub Copilot
powered by Anthropic Claude Opus 4.7, Anthropic Claude)
were used for prose drafting, consistency-checking, and
structural re-organisation; all conjectures, theorem
statements, and editorial decisions remain the author's.
All numerical claims trace to exact computations stored
under \texttt{sessions/2026-05-\{01,02\}/} on the SIARC
bridge with AEAL claim entries in \texttt{claims.jsonl}.

% =====================================================================
\bibliographystyle{plain}
\bibliography{annotated_bibliography}

\end{document}
